\section*{ВВЕДЕНИЕ}
\addcontentsline{toc}{section}{ВВЕДЕНИЕ}

В современном мире, где технологии играют ключевую роль в повседневной жизни, эргономика и удобство использования устройств становятся всё более важными аспектами. Одним из таких устройств является клавиатура, которая является неотъемлемой частью взаимодействия человека с компьютером. Однако традиционные клавиатуры часто не удовлетворяют потребности пользователей, особенно тех, кто проводит много времени за компьютером. Это привело к развитию специализированных решений, таких как сплит-клавиатуры, которые предлагают улучшенную эргономику и удобство использования.

Сплит-клавиатура Charybdis 4x6 Mk2 является одним из таких устройств, которое сочетает в себе уникальные особенности, такие как разделение на две части, трекбол и возможность горячей замены переключателей. Эти характеристики делают её привлекательной для пользователей, стремящихся к повышению комфорта и эффективности работы. Однако, несмотря на все преимущества, использование такой клавиатуры требует специализированного программного обеспечения, которое позволит максимально раскрыть её потенциал.

Целью данной дипломной работы является разработка программного обеспечения для клавиатуры Charybdis 4x6 Mk2, которое обеспечит удобное и эффективное управление устройством, а также позволит пользователям настраивать его под свои индивидуальные потребности. В работе будут рассмотрены основные аспекты проектирования и реализации программного обеспечения, включая выбор технологий, разработку архитектуры и тестирование готового продукта.

Развитие технологий и увеличение времени, проводимого за компьютером, привели к росту интереса к эргономичным решениям. Сплит-клавиатуры, такие как Charybdis 4x6 Mk2, предлагают улучшенную эргономику, что снижает нагрузку на руки и уменьшает риск возникновения заболеваний опорно-двигательного аппарата. Однако, для полноценного использования таких устройств необходимо специализированное программное обеспечение, которое позволит пользователям настраивать клавиатуру под свои индивидуальные потребности.

\emph{Цель настоящей работы} -- создание программного обеспечения для клавиатуры Charybdis 4x6 Mk2, которое обеспечит удобное и эффективное управление устройством. Для достижения этой цели необходимо решить следующие задачи:

\begin{enumerate}
\item Анализ существующих решений: изучение рынка программного обеспечения для сплит-клавиатур и выявление их преимуществ и недостатков.
\item Проектирование архитектуры: разработка архитектуры программного обеспечения, которая будет учитывать особенности клавиатуры и потребности пользователей.
\item Выбор технологий: определение наиболее подходящих технологий и инструментов для реализации программного обеспечения.
\item Реализация программного обеспечения: написание кода и создание пользовательского интерфейса.
\item Тестирование и отладка: проведение тестирования программного обеспечения и устранение выявленных ошибок.
\item Документирование: создание документации, которая поможет пользователям и разработчикам в использовании и поддержке программного обеспечения.
\end{enumerate}

\emph{Структура и объем работы.} Отчет состоит из введения, 4 разделов основной части, заключения, списка использованных источников, 2 приложений. Текст выпускной квалификационной работы равен \formbytotal{lastpage}{страниц}{е}{ам}{ам}.

\emph{Во введении} сформулирована цель работы, поставлены задачи разработки, описана структура работы, приведено краткое содержание каждого из разделов.

\emph{В первом разделе} на стадии описания технической характеристики предметной области приводится сбор информации о сплит-клавиатурах, а также о их преимуществах и недостатках.

\emph{Во втором разделе} на стадии технического задания приводятся требования к разрабатываемому программному обеспечению.

\emph{В третьем разделе} на стадии технического проектирования представлены проектные решения для программного обеспечения.

\emph{В четвертом разделе} приводится список классов и их методов, использованных при разработке программного обеспечения, производится тестирование разработанного программного обеспечения.

В заключении излагаются основные результаты работы, полученные в ходе разработки.

В приложении А представлен графический материал.
В приложении Б представлены фрагменты исходного кода. 
